\section{Introduction}
\label{sec:introduction}

Grading mathematical problem solutions is a cognitively demanding task that
requires not merely checking the final answer but tracing the logical
dependencies between intermediate steps, identifying where errors originate,
and producing a coherent written critique. In educational settings, expert
graders naturally construct an implicit \emph{reasoning graph}---a directed
acyclic structure linking premises to conclusions---when evaluating student
work. Recent advances in reinforcement learning with verifiable rewards
(RLVR)~\citep{deepseekr1,shao2024deepseekmath} have shown that language
models can be fine-tuned to solve and evaluate mathematical problems by
receiving binary correctness signals. However, these outcome-only rewards
leave the \emph{process} of reasoning entirely unsupervised, leading to
several well-documented failure modes.

\paragraph{Limitations of outcome-only RLVR.}
We identify three systematic shortcomings of current RLVR approaches when
applied to mathematical critique:

\begin{enumerate}[leftmargin=*,itemsep=2pt]
    \item \textbf{Outcome hacking.} Models learn to predict the correct
    verdict (e.g., ``the solution is wrong at step~3'') without
    constructing a logically sound justification. The reward signal
    provides no gradient for improving the quality of the reasoning
    trace itself.

    \item \textbf{Format dependence.} Because the verifier checks only
    the final structured output, models become brittle to minor
    formatting variations in the input solution, mistaking stylistic
    differences for logical errors.

    \item \textbf{Long-chain redundancy.} Without any length or
    coherence penalty, RLVR-trained models tend to produce verbose,
    repetitive reasoning chains. This wastes inference budget and often
    \emph{degrades} accuracy by introducing self-contradictions in
    longer outputs.
\end{enumerate}

\paragraph{TopoPRM: topology-aware process rewards.}
To address these limitations, we propose \textbf{TopoPRM}, a composite
reward framework that supplements the outcome reward with two
topology-aware process signals. Given a model-generated critique, we first
extract the implicit reasoning DAG by parsing step-level dependency
annotations. A \emph{topological structure reward} $R_{\mathrm{topo}}$
then evaluates whether this DAG satisfies basic structural properties
(acyclicity, connectivity, absence of orphan nodes). Concurrently, a
\emph{continuity reward} $R_{\mathrm{cont}}$ measures the semantic
coherence between consecutive reasoning steps using embedding similarity,
penalising abrupt logical jumps. Together with a lightweight format reward,
these signals compose a fully verifiable, code-computed reward that
requires no learned reward model.

We integrate TopoPRM into the Group Relative Policy Optimisation
(GRPO)~\citep{shao2024deepseekmath} framework and additionally propose a
\emph{reasoning compression} stage: after GRPO training, we distil the
improved reasoning behaviour from the 14B teacher into smaller 7B and 1.5B
student models via reverse-KL divergence minimisation, transferring both
accuracy and conciseness.

\paragraph{Contributions.}
Our main contributions are as follows:
\begin{enumerate}[leftmargin=*,itemsep=2pt]
    \item We propose TopoPRM, the first topology-aware verifiable process
    reward for RLVR that evaluates the structural quality of reasoning
    graphs without requiring a learned reward model.

    \item We design a composite GRPO reward combining outcome, topological
    structure, continuity, and format signals, all computed
    programmatically from model outputs.

    \item We introduce a reverse-KL distillation pipeline that compresses
    the reasoning improvements from a 14B teacher into 7B and 1.5B
    students with minimal accuracy loss.

    \item We conduct extensive experiments on MATH500, GSM8K, CMATH, and
    GaoKao, demonstrating that TopoPRM improves critique accuracy by
    \textbf{TODO}\% while reducing average reasoning length by
    \textbf{TODO}\% compared to outcome-only GRPO.
\end{enumerate}
